The same table shows the price of anchoring, and it is the one \cref{thm:contraction} predicts.
When no proposal can reach a mode (high noise, high dimension), $M_t$ is huge and the contraction factor $1-\epsilon_K$ is close to one. The target then follows the buffer: AEM keeps the modes and mode weights the buffer has, and it does not discover new ones.
Started from scratch, AEM inherits whichever modes the first, essentially random, buffers found. On GMM-10 it covers $5$ of $10$ modes (mode TV $0.65$), against $8$ for iDEM (mode TV $0.21$), and its $\mathcal{W}_2$ and MMD are worse accordingly. On MW-32 it occupies the minor wells less evenly than iDEM.
iDEM explores better precisely \emph{because} its target is biased. When SNIS collapses onto the proposal closest to a mode, the resulting field points every noised point towards its nearest low-energy region, which spreads mass across basins whatever the buffer contains.

Exploring first and anchoring afterwards recovers most of iDEM's coverage. On GMM-10 the explore$\to$anchor run keeps all 8 modes iDEM found and still halves $\mathcal{E}$-$\mathcal{W}_1$. It does not recover iDEM's mode-weight accuracy (mode TV $0.35$ vs.\ $0.21$).
Replay proposals were designed to correct mode weights through the energy (\cref{sec:replay}). In our runs they gave the best variance ratio and $\mathcal{W}_2$ on MW-32, but they did not close the mode-weight gap. Used \emph{without} an anchor from the start, they collapse onto the buffer's modes even faster (2 of 10 modes on GMM-10 after 10 iterations; \cref{app:extra}).

The overall picture is consistent. Anchoring makes the target unbiased exactly where the anchor is informative, and it corrects iDEM's systematic within-mode distortion. Mode discovery is a separate problem, and the bias that hurts iDEM's within-mode accuracy is what helps its exploration.
A sampler that wants both needs an exploration mechanism that does not rely on biased targets, such as tempering or annealing across noise levels. Anchoring can be combined with such mechanisms.
